% 木星（综述）
% license CCBYSA3
% type Wiki

本文根据 CC-BY-SA 协议转载翻译自维基百科\href{https://en.wikipedia.org/wiki/Jupiter}{相关文章}。


木星是距离太阳第五颗行星，也是太阳系中体积最大的行星。它是一颗气态巨行星，其质量几乎是太阳系中所有其他行星总和的 2.5 倍，但仍不到太阳质量的千分之一。它的直径是地球的 11 倍，是太阳的十分之一。木星以 5.20 天文单位（778.5 吉米）的平均距离绕太阳运行，公转周期为 11.86 年。在地球夜空中，木星是仅次于月球和金星的第三亮的天然天体，自史前时代以来便为人类所观察。它的名字源自古罗马宗教中至高神朱庇特（Jupiter）。

木星是太阳行星中最先形成的一颗，其在太阳系原始阶段向内迁移的过程深刻影响了其他行星的形成历史。木星的大气按质量计由 76\% 的氢和 24\% 的氦组成，内部密度更高。它还含有微量的碳、氧、硫、氖等元素，以及氨、水蒸气、磷化氢、硫化氢和烃类等化合物。木星的氦含量约为太阳的 80\%，与土星的成分相似。

木星外层大气被分为一系列纬向云带，这些云带在交界处产生湍流和风暴；其中最显著的结果就是“大红斑”——自 1831 年以来被记录的一场巨型风暴。由于木星自转极快（约 10 小时自转一周），使其呈现为一颗扁球体；与两极相比，赤道部位存在约 6.5\% \(^\text{[e]}\)的明显隆起。科学家认为，木星内部结构包括一层由流体金属氢组成的外地幔，以及由更高密度物质构成的弥散内核。木星内部持续的收缩过程释放出比其从太阳接收的还要多的热量。木星的磁场是太阳系中最强且连续范围第二大的结构，由流体金属氢内部的涡电流产生。太阳风与其磁层相互作用，使磁层向外延伸并影响木星的运行环境。

至少有 97 颗卫星绕木星运行；其中四颗最大的卫星——木卫一（Io）、木卫二（Europa）、木卫三（Ganymede）和木卫四（Callisto）——位于其磁层内，用普通双筒望远镜即可观测。四者中最大的木卫三比行星水星还要大。木星周围还存在一个暗淡的行星环系统。木星环主要由尘埃组成，分为三部分：内侧呈光环状尘粒环的“晕环”（halo）、相对明亮的主环（main ring），以及外侧的薄雾状“薄环”（gossamer ring）。这些行星环在可见光和近红外光下呈现红色。木星环系统的年龄尚不清楚，可能追溯至木星形成时期。自 1973 年以来，已有九个机器人探测器访问过木星：包括七次飞掠任务和两艘专门的轨道探测器（另有两艘在途中）。在其他行星系统中也已发现类似木星的系外行星。
\subsection{名称与符号}.
Jupiter 在古希腊和古罗马文明中都以神系中的最高神命名：对希腊人而言是宙斯（Zeus），对罗马人而言是朱庇特（Jupiter）\(^\text{[20]}\)。国际天文学联合会于 1976 年正式采用 Jupiter 作为该行星的名称，并自此将新发现的卫星以这位神祇的情人、宠爱者及后代命名\(^\text{[21]}\)。木星的行星符号 ♃ 源自带有水平笔画的希腊字母 zeta（⟨Ƶ⟩），作为 Zeus 的缩写\(^\text{[22][23]}\)。在拉丁语中，Iovis 是 Iuppiter（即 Jupiter）的属格形式，并与 Zeus（“天空之父”）的词源相关。英语中的对应形式 Jove 被认为在 14 世纪起作为该行星的诗性名称出现\(^\text{[24]}\)。Jovian 是 Jupiter 的形容词形式。较早的形容词 jovial，被中古时代的占星家使用，后来演变为“愉快的”或“欢乐的”含义，此情绪被认为与木星在占星学中的影响相关\(^\text{[25]}\)。原始希腊神祇 Zeus 还提供了词根 zeno-，用于构成一些与木星相关的词语，例如 zenography\(^\text{[f]}\)。
\subsection{形成与迁移}
木星被认为是太阳系中最古老的行星，其形成时间仅在太阳之后一百万年，约早于地球五千万年\(^\text{[26]}\)。当前的太阳系形成模型表明，木星形成于雪线处或其以外的位置：即距离早期太阳足够远、温度足够低，使水等挥发物能够凝结成固体的区域\(^\text{[27]}\)。木星首先形成固体核心，随后吸积其气态大气层。因此，该行星必须在太阳星云完全消散之前形成\(^\text{[28]}\)。在形成过程中，木星的质量逐渐增加，直到达到地球质量的 20 倍，其中约一半由硅酸盐、冰以及其他重元素成分构成\(^\text{[26]}\)。当原始木星增大超过 50 个地球质量时，它在太阳星云中开辟出一个空隙\(^\text{[26]}\)。此后，这颗不断成长的行星在 300–400 万年内达到最终质量\(^\text{[26][28]}\)。

根据“Grand tack 假说”，木星最初在距太阳约 3.5 AU（5.20 亿公里；3.30 亿英里）的位置开始形成。随着年轻的木星吸积质量，它与环绕太阳的气体盘相互作用，并受到土星轨道共振的影响，从而向内迁移\(^\text{[27][29]}\)。这扰乱了数颗在更靠近太阳轨道上运行的超级地球，使它们发生破坏性的碰撞\(^\text{[30]}\)。随后，土星开始以比木星更快的速度向内迁移，直到两者在距太阳约 1.5 AU（2.20 亿公里；1.40 亿英里）处被俘获进入 3:2 平均运动共振\(^\text{[31]}\)。这一事件改变了迁移方向，使它们开始远离太阳并迁出内太阳系，移动至如今的位置\(^\text{[30]}\)。上述全过程持续约 300–600 万年，而木星最终阶段的迁移历经数十万年完成\(^\text{[29][32]}\)。木星从内太阳系的迁移最终使包括地球在内的内行星得以从残余碎片中形成\(^\text{[33]}\)。

Grand tack 假说仍存在多个未解决的问题。其推导出的类地行星形成时间尺度似乎与测量得到的元素组成不一致\(^\text{[34]}\)。如果木星曾穿越太阳星云迁移，它最终可能会定居在更靠近太阳的轨道上\(^\text{[35]}\)。一些太阳系形成的竞争模型预测木星的形成轨道性质与今日行星的位置非常接近\(^\text{[28]}\)。其他模型则预测木星形成于更为外侧的位置，例如 18 AU（27 亿公里；17 亿英里）\(^\text{[36][37]}\)。

根据 Nice 模型，在太阳系历史最初的 6 亿年间，原始柯伊伯带天体的内落导致木星和土星从它们最初的位置迁移至 1:2 共振，这使土星移入更高轨道，从而扰乱了天王星和海王星的轨道，消耗了柯伊伯带，并触发了晚期重轰炸\(^\text{[38]}\)。

根据 Jumping-Jupiter 情景，木星在早期太阳系中的迁移可能导致第五颗气态巨行星被抛出。该假说认为，在其轨道迁移过程中，木星的引力扰动破坏了其他气态巨行星的轨道，可能使其中一颗完全被逐出太阳系。此类事件的动力学将极大改变太阳系的形成与配置，最终只留下如今人类所见的四颗气态巨行星\(^\text{[39]}\)。

基于木星的组成，有研究者提出其最初形成于分子氮（N\(_2\)）雪线之外，即距太阳 20–30 AU（30–45 亿公里；19–28 亿英里）的位置，甚至可能形成于氩雪线之外，后者可能远至 40 AU（60 亿公里；37 亿英里）\(^\text{[40][41]}\)。如果木星形成于这些极远的位置之一，则它将在大约 70 万年的时间尺度上向内迁移至目前的位置\(^\text{[36][37]}\)，这一过程发生于行星开始形成后约 200–300 万年的时代。在该模型中，土星、天王星和海王星的形成位置将比木星更远，而土星亦会向内迁移\(^\text{[36]}\)。
\subsection{物理特征}
木星是一颗气态巨行星，这意味着其化学组成主要为氢和氦。这些材料在行星地质学中被分类为“气体”，这一术语并不指示其物质状态。木星是太阳系中最大的行星，其赤道直径为 142,984 公里（88,846 英里），体积为地球的 1,321 倍\(^\text{[3][42]}\)。其平均密度为 1.326 g/cm\(^3\)\(^\text{[g]}\)，低于四颗类地行星的密度\(^\text{[44][45]}\)。
\subsubsection{成分}
木星大气按质量计约由 76\% 的氢和 24\% 的氦组成。按体积计，上层大气约含 90\% 的氢和 10\% 的氦，这一较低比例是因为单个氦原子比该大气层中形成的氢分子更为致密\(^\text{[46]}\)。大气含有微量的碳、氧、硫和氖\(^\text{[47]}\)，以及氨、水蒸气、磷化氢、硫化氢和甲烷、乙烷、苯等烃类化合物\(^\text{[48]}\)。其最外层包含冻结氨的晶体\(^\text{[49]}\)。行星内部密度更高，按质量计约由 71\% 的氢、24\% 的氦和 5\% 的其他元素组成\(^\text{[50][51]}\)。

大气中氢与氦的比例接近理论上的原始太阳星云成分\(^\text{[52]}\)。上层大气中的氖含量为每百万份质量中的 20 份，约为太阳丰度的十分之一\(^\text{[53]}\)。木星的氦含量约为太阳的 80\%，这是由于这些元素在行星内部深处以富含氦的液滴形式沉降的过程所致\(^\text{[54][55]}\)。

基于光谱学分析，土星的组成被认为与木星相似，但另外两颗巨行星天王星和海王星含有相对更少的氢和氦，而相对含有更多的次常见元素，包括氧、碳、氮和硫\(^\text{[56]}\)。这些行星被称为冰巨行星，因为在其形成期间，这些元素被认为以冰的形式并入其中；然而它们内部实际上可能含有极少量的冰\(^\text{[57]}\)。
\subsubsection{大小与质量}
\begin{figure}[ht]
\centering
\includegraphics[width=6cm]{./figures/82e5bc11e94dd9a1.png}
\caption{木星与地球及地球月球的尺寸比较} \label{fig_Jupite_1}
\end{figure}

